\paragraph{Problem.} Given a CT-like greyscale volume of a single object on an $N^3$ grid, estimate the
area of the object's true, continuous surface. The volume is degraded by the finite voxel size
(partial-volume averaging), by the scanner blur and by noise.

\paragraph{Objectives.} (1) Build shapes with reference areas accurate to better than $10^{-4}$, so that
the reference adds no error of its own. (2) Measure the error of voxel-face counting and three
marching-cubes variants across resolution and noise. (3) Train a 3D CNN and a PointNet to correct the
greyscale marching-cubes area, and compare them with a per-regime lookup table and a Random Forest.
(4) Test the corrections on shape families never seen in training. (5) State when a learned correction
is worth using.

\paragraph{Why this matters.} A 5\% error in surface area is large compared with the effects such
measurements are used to detect, for example differences between catalyst batches or the deviation of
a printed lattice from its design. Errors of this size are typical for small or noisy objects
(Section~\ref{sec:results}). A correction learned once on synthetic data is cheap to apply to every new
object.
